\section{Q3: a connectivity test from \texorpdfstring{$A$}{A}}\label{sec:test}

The test rests on two lemmas about walks and paths. \Cref{lem:cut} is Corollary~2.9.7 of Grinberg \cite[p.~39]{grinberg2025graph}; \cref{lem:reach} combines the walk-counting theorem with the bound that a path has at most $N-1$ edges (it has at most $N$ distinct vertices), which Grinberg states in Section~2.9.4 \cite[p.~40]{grinberg2025graph}. The proofs below are written out in full.

\begin{lemma}[Every walk contains a path]\label{lem:cut}
Let $W=u_0\cdots u_k$ be a walk from $u$ to $w$. Then there is a path from $u$ to $w$ of length at most $k$.
\end{lemma}
\begin{proof}
By strong induction on $k$ (strong, because $W'$ below can have any length from $0$ to $k-1$): for a walk of length $k$ we assume the statement for all walks of smaller length.

\emph{Base case $k=0$.} $W=(u)$ has one vertex, so it is a path of length $0$.

\emph{Inductive step, $k\ge1$.} If the vertices of $W$ are pairwise distinct, $W$ is itself a path of length $k$. Otherwise there are positions $s<t$ with $u_s=u_t$. Delete the vertices at positions $s+1,\dots,t$; this removes the second copy $u_t$ of the repeated vertex and everything between the two copies, and gives $W'=u_0\cdots u_su_{t+1}\cdots u_k$, and $W'=u_0\cdots u_s$ if $t=k$. $W'$ is a walk, because its only new consecutive pair is $u_s,u_{t+1}$, and $u_s=u_t$ is adjacent to $u_{t+1}$. It starts at $u_0=u$, ends at $u_k=w$ (if $t=k$ it ends at $u_s=u_t=u_k$), and has length $k-(t-s)<k$. The induction hypothesis, applied to the shorter walk $W'$, gives a path from $u$ to $w$ of length at most $k-(t-s)\le k$.
\end{proof}

\begin{lemma}[Reachability from the powers]\label{lem:reach}
Let $i\ne j$. There is a path from $v_i$ to $v_j$ if and only if $(A^k)_{ij}>0$ for some $k$ with $1\le k\le N-1$.
\end{lemma}
\begin{proof}
($\Rightarrow$) Let a path from $v_i$ to $v_j$ have length $k$. As $i\ne j$, $k\ge1$ (a path of length $0$ has a single vertex, so it cannot join different vertices). A path of length $k$ has $k+1$ pairwise distinct vertices among the $N$ vertices of the graph, so $k+1\le N$, that is $k\le N-1$. A path is a walk, so by \cref{thm:counting} $(A^k)_{ij}\ge1$.

($\Leftarrow$) If $(A^k)_{ij}>0$ with $1\le k\le N-1$, then by \cref{thm:counting} a walk of length $k$ joins $v_i$ to $v_j$, and by \cref{lem:cut} a path joins them.
\end{proof}

\begin{theorem}[Sum test]\label{thm:sum}
Let $N\ge2$ and $B=A+A^2+\dots+A^{N-1}$. Then $G$ is connected if and only if $B_{ij}>0$ for all $i\ne j$.
\end{theorem}
For $N=2$ the sum has the single term $B=A$.
For $N=1$ there is no pair of different vertices and the graph is connected. \Cref{thm:sum} is Theorem~2 of Paramadevan and Sotheeswaran \cite[p.~17]{paramadevan2021adjacency}.
\begin{proof}
Every $A^k$ has non-negative entries, so $B_{ij}=\sum_{k=1}^{N-1}(A^k)_{ij}$ is a sum of non-negative terms, and it is positive exactly when at least one term is positive. For $i\ne j$ this holds exactly when a path joins $v_i$ and $v_j$, by \cref{lem:reach}. So $B_{ij}>0$ for all $i\ne j$ exactly when every pair of different vertices is joined by a path, which is \cref{def:connected}.
\end{proof}
The cutoff $N-1$ is the largest length a path can have (proof of \cref{lem:reach}). The entries of $B$ are walk counts, not path counts; only whether they are zero is used. The test can also be written as ``$(I+A)^{N-1}$ has no zero entry'' \cite[Proposition~3, p.~2]{gotti18211}\cite[Theorem~2, p.~1]{bai2025adjacency}. Indeed $(I+A)^{N-1}=\sum_{k=0}^{N-1}\binom{N-1}{k}A^k$ because $I$ and $A$ commute; every coefficient is at least $1$, the term $k=0$ makes the diagonal positive, and off the diagonal only the terms $k\ge1$ contribute, each a positive multiple of $(A^k)_{ij}$ (every coefficient is at least $1$), so the entry is positive exactly when some $(A^k)_{ij}$ with $1\le k\le N-1$ is positive, that is, exactly when $B_{ij}>0$.
